\chapter{LDH Isoenzymes}

\section*{What Is This Test?}
The LDH isoenzyme test separates and quantifies the five isoenzymes of lactate dehydrogenase (LDH-1 through LDH-5) to determine the tissue source of elevated total LDH. LDH is a tetrameric enzyme composed of four polypeptide subunits of two types: H (heart) and M (muscle), encoded by the LDHA and LDHB genes on chromosomes 11 and 12, respectively. The five isoenzymes, in order of electrophoretic mobility from fastest (most anodal) to slowest (most cathodal), are: LDH-1 (H4, HHHH)---found primarily in cardiac muscle, red blood cells, and renal cortex; LDH-2 (H3M, HHHM)---found in cardiac muscle, RBCs, and reticuloendothelial system; LDH-3 (H2M2, HHMM)---found in lungs, lymphocytes, pancreas, and spleen; LDH-4 (H1M3, HMMM)---found in liver, skeletal muscle, and kidneys; LDH-5 (M4, MMMM)---found in liver and skeletal muscle. In normal serum, the isoenzyme pattern is: LDH-2 > LDH-1 > LDH-3 > LDH-4 > LDH-5. A ``flipped'' LDH pattern (LDH-1 > LDH-2) is characteristic of myocardial infarction and hemolysis (both RBCs), because the heart and RBCs contain predominantly LDH-1. Predominance of LDH-5 suggests liver or skeletal muscle injury. The LDH isoenzyme test has been largely supplanted by more specific markers (troponin for MI, ALT/AST for liver, CK for muscle) but remains useful when the source of LDH elevation is unclear, or when evaluating certain malignancies (germ cell tumors, where LDH-1 elevation is seen) \cite{tietz2023textbook}.

\section*{Alternative Names}
Lactate Dehydrogenase Isoenzymes; LDH Fractionation; LDH Isoenzyme Electrophoresis; LD Isoenzymes; LDH-1, LDH-2, LDH-3, LDH-4, LDH-5

\section*{Principle}
LDH isoenzymes are separated by agarose gel electrophoresis. The serum sample is applied to the gel and subjected to an electric field; the five isoenzymes migrate at different rates based on their net charge (LDH-1, being most negatively charged, migrates fastest toward the anode). After electrophoresis, the gel is overlaid with a substrate solution containing lactate, NAD\textsuperscript{+}, phenazine methosulfate, and nitrotetrazolium blue (NTB). LDH activity in each band reduces NAD\textsuperscript{+} to NADH, which reduces NTB to a purple/blue formazan precipitate. The intensity of staining in each band is quantified by densitometry. Results are reported as a percentage of total LDH for each isoenzyme and as absolute activity (IU/L) \cite{panteghini2023serum}. Heat stability testing can also be used: LDH-1 is heat-stable at 65 degrees C for 30 minutes; LDH-5 is heat-labile. Chemical inhibition (urea inhibits LDH-5; oxalate inhibits LDH-1) can provide additional confirmation. A ``flipped LDH pattern'' is defined as LDH-1 greater than LDH-2 (normally LDH-2 is 30--40\% and LDH-1 is 17--27\%) \cite{bais1994approved}. This test is not performed by hematology analyzers.

\section*{History}
LDH isoenzymes were first separated by Vesell and Bearn in 1957 using starch gel electrophoresis \cite{vesell1957localization}. The diagnostic utility of the flipped LDH pattern (LDH-1 > LDH-2) for myocardial infarction was established in the 1960s \cite{lott1980serum}. Before troponin, LDH isoenzymes were a standard part of the MI diagnostic panel because LDH-1 remains elevated long after CK-MB has normalized (10--14 days vs. 2--3 days), allowing late MI diagnosis. The clinical role of LDH isoenzymes declined with the advent of troponin in the 1990s. Current clinical applications are primarily for hematologic disorders (hemolysis) and select malignancies.

\section*{Why This Test Is Ordered}
\begin{itemize}
  \item Confirmation of hemolysis as the cause of elevated total LDH: flipped LDH pattern (LDH-1 > LDH-2) with elevated total LDH supports hemolysis or myocardial injury
  \item Differentiation of hemolysis (LDH-1 predominant) from liver disease (LDH-5 predominant) when total LDH is elevated and clinical context is ambiguous
  \item Evaluation of disseminated germ cell tumors (testicular cancer): elevated LDH-1 is associated with seminoma and non-seminomatous germ cell tumors
  \item Investigation of unexplained, persistent elevation of total LDH without obvious cause: isoenzyme pattern may point to tissue source \cite{wolf1986interpretation}
  \item Adjunctive diagnosis of myocardial infarction when troponin is unavailable or equivocal and presentation is late (3--7 days after symptoms): flipped LDH ratio persists for 10--14 days \cite{lott1980serum}
  \item Detection of renal infarction: LDH-1 and LDH-2 are markedly elevated due to high LDH-1 content in renal cortex
\end{itemize}

\section*{Primary vs Secondary Test}
LDH isoenzymes are a secondary, specialized test. They are rarely ordered in modern clinical practice due to the availability of more specific organ markers (troponin, ALT/AST, CK). The test is reserved for specific clinical scenarios where tissue source identification of LDH is needed.

\section*{How This Test Is Performed}
A healthcare professional draws blood from a vein into a serum separator tube (gold-top). The sample must be completely free of hemolysis, as red blood cell LDH-1 contamination will produce a falsely abnormal isoenzyme pattern. The serum is separated promptly and the isoenzyme electrophoresis is performed in a specialized laboratory. Results are available within 2--5 days (this is a send-out test in many hospitals). Some institutions use immunochemical methods to specifically quantify LDH-1 (anti-M subunit inhibition) without full electrophoresis, providing faster results.

\section*{What Preparation Is Needed}
No fasting is required. The critical requirement is avoidance of hemolysis during phlebotomy. Trauma from difficult venipuncture, small needles, excessive suction, and shaking the tube must be avoided. Recent exercise, surgery, or intramuscular injections can alter the isoenzyme pattern.

\section*{Contraindications and Precautions}
No absolute contraindications. Critical consideration: hemolyzed specimens CANNOT be used because RBC LDH-1 contamination creates a false flipped LDH pattern. If hemolysis is present, the test is uninterpretable and must be redrawn.

\section*{Risks}
Minimal risk from venipuncture.

\section*{What Happens After the Test}
Results available in 2--5 days. The normal isoenzyme distribution is: LDH-2 > LDH-1 > LDH-3 > LDH-4 > LDH-5. Flipped LDH (LDH-1 > LDH-2): hemolysis (most common), MI (historical), renal infarction, germ cell tumors (testicular cancer). Elevated LDH-3: pulmonary embolism/infarction, pancreatitis, lymphoma. Elevated LDH-4 and LDH-5: liver disease (hepatitis, cirrhosis, hepatic congestion, metastases), skeletal muscle injury (rhabdomyolysis, myositis, trauma, exercise). This test has limited clinical utility in modern practice due to superior organ-specific markers.

\section*{Understanding Results}

\subsection*{Below Normal}
\begin{itemize}
  \item Not clinically significant. Decreased individual isoenzyme fractions have no defined clinical interpretation.
\end{itemize}

\subsection*{Above Normal}
\begin{itemize}
  \item \textbf{Flipped LDH pattern (LDH-1 > LDH-2):} LDH-1 >27\%, LDH-1/LDH-2 ratio >1.0. Causes: (a) Hemolytic anemia (most common current reason for ordering): intravascular hemolysis from any cause. (b) Myocardial infarction (historical): appears 12--24 hours after MI, persists for 10--14 days. (c) Renal infarction: renal cortex is rich in LDH-1. (d) Germ cell tumors (testicular cancer): seminoma and non-seminomatous tumors. (e) Megaloblastic anemia: ineffective erythropoiesis releases LDH-1 from destroyed marrow erythroblasts.
  \item \textbf{Elevated LDH-3:} Pulmonary embolism/infarction, acute pancreatitis, lymphoma, acute leukemia, lymphocytosis. Lung and lymphoid tissues are rich in LDH-3.
  \item \textbf{Elevated LDH-4 and LDH-5:} (a) Liver disease: acute hepatitis (viral, toxic, alcoholic), cirrhosis, hepatic congestion (right heart failure, Budd-Chiari syndrome), hepatic metastases. (b) Skeletal muscle injury: rhabdomyolysis, polymyositis/dermatomyositis, trauma, strenuous exercise. (c) Hypothyroidism. (d) Malignancy (many solid tumors show LDH-5 predominance).
  \item \textbf{Pan-isoenzyme elevation:} Hemolysis (RBCs contain LDH-1, -2, and -3), widespread malignancy, sepsis, shock, and critical illness.
\end{itemize}

\section*{Normal Results}
\begin{itemize}
  \item \textbf{LDH-1 (H4, heart/RBC):} 17--27\% of total (normally less than LDH-2). Absolute: 20--75 IU/L \cite{wallach2022interpretation}
  \item \textbf{LDH-2 (H3M):} 27--37\% of total (normally the highest fraction). Absolute: 30--100 IU/L
  \item \textbf{LDH-3 (H2M2):} 18--25\% of total. Absolute: 20--70 IU/L
  \item \textbf{LDH-4 (H1M3):} 8--16\% of total. Absolute: 10--45 IU/L
  \item \textbf{LDH-5 (M4, liver/muscle):} 6--16\% of total. Absolute: 7--45 IU/L
  \item \textbf{Normal pattern:} LDH-2 > LDH-1 > LDH-3 > LDH-4 > LDH-5
  \item \textbf{Flipped pattern:} LDH-1 > LDH-2 (LDH-1/LDH-2 ratio >1.0)
\end{itemize}

\section*{Follow-up Tests}
\begin{itemize}
  \item \textbf{Corresponding organ-specific markers:} ALT/AST/ALP/GGT for liver; troponin for heart; CK for muscle.
  \item \textbf{CBC, reticulocyte count, haptoglobin, peripheral blood smear:} For hemolysis workup if LDH-1 elevated.
  \item \textbf{Abdominal imaging (ultrasound, CT):} For hepatic or renal pathology.
  \item \textbf{Tumor markers:} AFP, beta-hCG for germ cell tumors; LDH is itself elevated.
  \item \textbf{Pulmonary imaging (CTPA):} For suspected PE.
  \item \textbf{LDH in other body fluids:} Pleural fluid LDH for Light's criteria (exudate vs transudate); CSF LDH for CNS malignancy or infection.
\end{itemize}

\section*{References}
\bibliographystyle{unsrtnat}
\bibliography{references}

\clearpage
\section*{Regulatory Status and Authorization Date}
Regulatory status applies to the specific assay, instrument, manufacturer, and intended use rather than to the test name alone. No single approval date can be assigned to this test category. For a commercial assay, verify the current product in the relevant regulator's database: the U.S. Food and Drug Administration (FDA) clearance, approval, or authorization; India's Central Drugs Standard Control Organisation (CDSCO) licence or permission; the European Union In Vitro Diagnostic Regulation (EU IVDR) CE certificate issued through the applicable conformity-assessment route; or the United Kingdom Medicines and Healthcare products Regulatory Agency (MHRA) registration and UKCA/CE status. Laboratory-developed tests may not have an individual FDA, CDSCO, EU IVDR, or MHRA approval date.
\clearpage
